% Generated draft: requires mathematical and source review.
\documentclass[11pt]{article}
\usepackage[margin=25mm]{geometry}
\usepackage{fontspec}
\setmainfont{Latin Modern Roman}
\usepackage{amsmath,amssymb,mathtools,unicode-math}
\setmathfont{Latin Modern Math}
\usepackage{enumitem,needspace,xurl,hyperref}
\setlength{\parindent}{0pt}
\setlength{\parskip}{5pt}
\setlength{\emergencystretch}{3em}
\newcommand{\N}{\mathbb{N}}
\newcommand{\Z}{\mathbb{Z}}
\newcommand{\Q}{\mathbb{Q}}
\newcommand{\R}{\mathbb{R}}
\newcommand{\C}{\mathbb{C}}
\newcommand{\F}{\mathbb{F}}
\newcommand{\A}{\mathbb{A}}
\newcommand{\PP}{\mathbb{P}}
\newcommand{\E}{\mathbb{E}}
\newcommand{\eps}{\varepsilon}
\newcommand{\dd}{\,\mathrm{d}}
\begin{document}
\section*{UnsolvedMath ID: 10400223}
\textbf{Source problem code:} catalogue entry
\subsection{Definitions and mathematical statement}
J. Roberts asks to extend Kuperberg's work on webs. In the usual representation-theoretic usage, a web is a planar graph with suitable edge labels, orientations and vertex types, representing an invariant tensor or intertwining map; relations among webs encode relations among these maps.

The target is to construct a generalization of the particular web theory intended in the source, with a specified new domain and a comparison recovering the existing theory in its original cases. The excerpt does not identify Kuperberg's paper, the original groups or categories, the intended enlarged domain, the parameter or coefficient conventions, or which properties the extension must retain. In particular, no fixed higher-rank presentation, basis theorem or completeness assertion is printed. Those would be possible precise versions only after the omitted specification is recovered.
\subsection{Short English statement}
Develop an extension of the intended Kuperberg web theory, specifying the new domain and how it recovers the original theory; the excerpt does not fix those choices.
\subsection{Sources}
\begin{itemize}
\item J. Roberts, attributed Problem 12.18 in T. Ohtsuki, \emph{Problems on invariants of knots and 3-manifolds} (2002), supplied PDF page 167. \url{https://msp.org/gtm/2002/04/gtm-2002-04-024s.pdf}. Recorded access 2026-07-29; extraction pdf-font-structure. Kuperberg's work is named without a paper title or bibliography entry.
\item Local UnsolvedMath registry, \texttt{data/unsolvedmath/problems.json}, permanent ID \(10400223\), exact source code \texttt{catalogue entry}. The supplied OPEN-TRIAGE summary says that various authors extended web theory to higher-rank groups, but supplies no names, papers or exact scope.
\end{itemize}
\end{document}
